\documentclass{jsarticle}
\usepackage[dvipdfmx,final]{graphicx}
\usepackage{amsmath}
\usepackage{amssymb}
\begin{document}
\title{}
\author{}
\date{}
\maketitle

   $$ e^{f} + e^{-f} \ge e^{f} - e^{-f} $$
   $$ \bigoplus{(i \hbar^{\nabla})}^{\oplus L} = e^{-f} $$
   $$ {L^{\nabla}}^{\oplus L} - {{L^{\nabla}}^{\oplus L}}^{-1}
   \le {L^{\nabla}}^{\oplus L} + {{L^{\nabla}}^{\oplus L}}^{-1} \le 
   e^{f} - e^{-f} \le e^{-f} + e^{f} $$
   $$ \left({L^{\nabla}}^{\oplus L} + {{L^{\nabla}}^{\oplus L}}^{-1}\right)^{df}
   = \sin \theta \ge {\sqrt{3} \over 2} $$ 
   $$ {d \over df}F + \int C dx_m \ge {d \over df}F - \int C dx_m $$
   $$ \ge {\sqrt{3} \over 2} $$
   オイラーの定数は、正常な細胞のヒッグス場のエネルギーからガンマ関数の
   大域的積分多様体であるオイラーの定数の多様体積分のエネルギーを
   差し引いたエネルギーであり、宇宙から異次元への角度の数値でもあり、
   このオイラーの定数のエネルギーバランスが、永遠の生命エネルギーの
   数値へのバランスであり、彩さんの過去のレビューからの導きでもあり、
   私の集大成であり、広島大学医学部からの導きであることは、
   私の論文からも、一目瞭然であり、富山大学への感謝である。
   $$ \beta(p,q) = \int||\sin^{2}\theta||d\tau $$
   $$ (\sin \theta \ge {\sqrt{3} \over 2}) $$
   の$ \theta \ge 60$度の範囲である。

  ベータ関数をガンマ関数へと渡すと、
   Gamma function sented for Beta function of inspiration with monitonicity
   of component with differential and integral being definitions.
   $$ \beta(p,q) = {{\Gamma(p)\Gamma(q)} \over {\Gamma(p + q)}} $$
   ここで、単体的微分と単体的積分を定義すると
   This area defined with different and integrate of component with
   global topology.
   $$ t = \Gamma(x) $$
   $$ T = \int \Gamma(x) dx $$
   これを使うのに、ベータ関数を単体的微分へと渡すために、
   for this system used for being is conbiniate 
   with beta function for componenet of deprivate eqaution.
   $$ \beta(p,q) = - \int {1 \over {t^2}}dt $$
   大域的微分変数へとするが、
   This point be for global differential variable exchanged,
   $$ T_m = \int \Gamma(x) dx_m $$
   この単体的積分を常微分へとやり直すと、
   This also point be for being retried from ordinary differential computation
   for component of integral manifold.
   $$ T^{'} = {t^{'} \over {{\log t}}}dt + C(\text{Cは積分定数}) $$
   これが、大域的微分多様体の微分変数と証明するためには、
   This exceed of proof being for being defined with deprivate variable of
   global differential formula, this successed for true is,
   $$ dx_m = {1 \over {\log x}}dx, dx_m = {(\log x^{-1})}^{'}  $$
   これより、単体的微分が大域的部分積分へと置換できて、
   Therefore, this exchanged from mononotocity deprivation from global
   parital integral manifold is able to, 
   $$ T^{'} = \int \Gamma(\gamma)^{'} dx_m $$

   微分幾何へと大域的積分と大域的微分多様体が、加群分解の共変微分へと
   書き換えられる。
   After all, these exchanged of monotonicity deprivation successed from
   being catastrophe of summativate of partial and assemble of deprivations
   for differential geometory of quantum level to global integral and
   differential manifold.
   $$ \bigoplus T^{\nabla}  = \int T dx_m, \delta(t) = t dx_m  $$


   this exluded of being conclution are which beta function
   evaluate with mononicity from ordinary differential equation
   be resulted from component of deprivation and integral expalanations.
   This cirtutation be resenbled to define with global topoloty of
   extention extern of deprivate and integral of manifolds
   estourced with quantum level of differential geometry be proofed with
   all of equation anbrabed from Euler product.
   これらをまとめると、ベータ関数を単体的関数として、
   これを常微分へと解を導くと、単体的微分と単体的積分
   へと定義できて、これより、大域的微分多様体と大域的積分多様体
   へとこの関数が存在していることを証明できるのを上の式たちから
   わかる。この単体的微分と単体的積分から微分幾何が書けることがわかる。



   虚数の虚数倍した値が超越数のx倍と同じとすると、
   超越数の$2 \pi m$倍が$n=1$でiとなると定義すると、
   次の式たちが導かれる。　
   Imaginary pole circlate with twigled of pole in step function
   from Naipia number of assembled from equalation of defined are
   escourted to be defined with next equations.


   These defined equation are climbate with idea of equation from
   Caltan of imaginary number of circulation.
   カルタンを超えているアイデアと数式たちでもある。

   $$ i^i = (\sqrt{i})^{\sqrt{i}} = e^{x \log x} $$
   $$ e^x = i, e^{2\pi m} = i^n, {d \over d x}e^{2\pi m} = i^n $$
   $$ e^{i \theta} = \cos \theta + i \sin \theta, e^{2 \pi m} = i^n $$
   $$ 2 \pi m = n, e = i \text{($e=i$となる$e^{i \theta} = \cos \theta 
   + i \sin \theta $
   の範囲で$2\pi m$ がある。)}$$
   $$ {d \over d i}i = i^i = e^{x \log x}, m={1 \over 2 \pi},l=2\pi r, 
   \pi = {l \over 2r} $$
   $$ ||ds^2|| = e^{-2 \pi T||\psi||}[\eta_{\mu\nu} + \bar{h}(x)]
   dx^{\mu}dx^{\nu} + T^2 d^2 \psi $$
   $$ e^x = i, e^{2\pi m} = i^n, e = i $$
   $$ {{[\eta_{\mu\nu} + \bar{h}(x)]dx^{\mu}dx^{\nu}}/i} = H \Psi
   = i \hbar \psi, H \Psi = {1 \over i}[H, \Psi] $$
   ハイゼンベルク方程式が$ AdS_5 $多様体の原子レベルの方程式も表せられる。 
   微分幾何の量子化の式は、
   Hisenbelg equation are represeted from AdS5 manifold to particle level
   equalation of quantum levle of differential geometry entranced.
   $$ \bigoplus {({i \hbar} ^{\nabla})}^{\oplus L} = e^{{i x \log x}^{
	   {e^{i x \log x}}^{'}}} $$
    $$ H\Psi = \bigoplus {H \Psi \over \nabla L} $$

   この式が、アカシックレコードの子どもの式である。２種類の表し方である。
   計算方法が、大域的微分多様体の求め方と同じである。

   移項して、
   それを展開していき、
    これらより、まとめた式が、
    $$ n!^{-1} = \bigoplus \nabla L $$
    これは、積分再発見法と同じであり、
    $$ n!^{-1}  = \int T^{{\mu\nu}^{'}} dV - \int T^{\mu\nu} dV $$
    これもカルーツァ・クライン空間と同じ式であり、

    $$ = \int f{(x)^{-1}} xf(x) - f(x) $$
    これに条件をつけると、
    $$ \exists x = a, 1 - \exists \int x dx, \Lambda = \int A dx $$
    これより、次の式になり、
    $$ \nabla f(x) = {1 \over n!} \int \int \cdots \int dx $$

    $$ \sum \int (\int \cdots A\int dz)/n! $$
    ウィッテン方程式により、
    $$ M^2 a^{(m)} = k a^{(m)} $$
    $$ ZM^3 = \int dAe^{ki/4\pi}\mathcal{L}M^3 $$
    $$ Z_M = \langle M_1 | M_2 \rangle $$
    ゼータ関数へといける。
    $$ e^{x} = \sum {x^n \over n!} (r = \infty) $$
    これらは、大域的微分多様体の部品たちになっている。
    $$ {d \over df}\int dx_m = {1 \over n!}\sum x^n $$


   Euler productは、
   $$ \int^{n+1}_{n}{dx \over x} < {1 \over n}, (n \ge 1) $$
   $$ 1 + {1 \over 2} + \cdots + {1 \over n} > \int^{n+1}_{1}{dx \over x}
   = \log(n+1), $$
   $$ 1 + {1 \over 2} + \cdots + {1 \over n} - \log x > \log{{(n+1)} 
   \over n} > 0, $$
   $$ {1 \over {n+1}} < \int^{n+1}_{0} {dx \over x} = \log(n+1) - \log n. $$
   $$ \lim_{n \to \infty} \left(1 + {1 \over 2} + \cdots + {1 \over n}
   - \log n \right) = C $$
   $$ = \int (\int {1 \over x^s}dx - \log x)\mathrm{dvol} = \int C dx_m $$
   Cの値は、$\text{C = 0.5772156$\cdots$} $である。

  オイラーの定数は、次の式で成り立っている。
  $$ C = \int{1 \over x^s}dx - \log x $$
  その式を単体積分をすると、
  $$ = \int {\left({\int {1 \over x^s}}dx - \log x\right)}\mathrm{dvol} $$
  ゼータ関数を多重積分すると、　
  $$ = \int \int {1 \over x^s}dx - \int {\log x} \mathrm{dvol} $$
  ここで、対数方程式は、多様体積分がガンマ関数を微分した方程式と
  同値より、
  $$ F = \Gamma = \int e^{-x}x^{1-t}dx $$
  $$ = {d \over df}\int \int {1 \over (y \log y)^{1 \over 2}}dx_m
  - \int e^{-x}x^{1 - t} \log x dx_m $$
  $$ \int x^{1 - t}e^{-x}dV = \int x^{1 - t}dm $$
  $$ \int x^{1 - t}e^{-x}dV = \int x^{1 - t}\mathrm{dvol} $$
  $$ f = \gamma = \Gamma^{'} = \int e^{-x}x^{1 - t} \log x dx $$
  $$ = {d \over d \gamma}\Gamma^{-1} - 
  \int e^{-x}x^{1 - t}\log x dx_m $$
  $$ = {d \over d \gamma}\Gamma^{-1} - \left(\gamma\right)^{\gamma^{'}} $$
  これらより、大域的微分方程式のヒッグス場方程式の逆三角関数の双曲多様体
  に対極する、双曲多様体の余弦定理と同値になる。
  $$ = e^{-f} - e^{f} $$
  $$ = 2 \cos (i x \log x) $$
  結局は、微分幾何の量子化がガンマ関数でもあり、その逆関数はゼータ関数
  でもあり、そして、オイラーの定数を多様体微分をすると、ヒッグス場の
  方程式は正弦定理の逆三角関数であり、このヒッグス場の
  方程式の逆三角関数が、余弦定理の
  双曲多様体として、オイラーの定数になる。
  オイラーの定数は、このガンマ関数から、無理数と証明される。

    カルーツァ・クライン空間の方程式は、
    $$ ds^2 = g_{\mu\nu}(x)dx^{\mu}dx^{\nu} + \psi^2(x)(dx^2 + \kappa^2
    A_{\mu}(x)dx^{\nu})^2 $$
    この式は、
    $$ ds^2 = -N(r)^2dt^2 + \psi^2(x)(dr^2 + r^2d\theta^2) $$
    $$ ds^2 = -dt^2 + r^{-8\pi Gm}(dr^2 + r^2d\theta^2) $$
    とまとまり、
    $$ dx^2 = g_{\mu\nu}(x)(g_{\mu\nu}(x)dx^2 - dxg_{\mu\nu}(x)) $$
    $$ dx = (g_{\mu\nu}(x)^2dx^2 - g_{\mu\nu}(x)dx g_{\mu\nu}(x))^{1 \over 2} $$
    と反重力と正規部分群の経路和となり、　
    $$ \pi(\chi, x) = i \pi (\chi, x)f(x) - f(x)  \pi (\chi, x) $$
    と基本群にまとまる。
    シュワルツシルト半径は、
    $$ ds^2 = -(1 - {r_s \over r})c^2 dt^2 + {1 \over {{1 - {r_s \over r}}}} 
    dr^2 + r^2 d \theta^2 + r^2 \sin^2 \theta d \psi^2 $$
    これは、カルーツァ・クライン空間と同型となる。
    一般相対性理論は、
    $$ R^{\mu\nu} + {1 \over 2}g_{ij}\Lambda = \kappa T^{\mu\nu} $$
    この多様体積分は、
    $$ \int \kappa T^{\mu\nu}\mathrm{dvol} = \int  \left(R^{\mu\nu} + 
    {1 \over 2}g_{ij}\Lambda\right)\mathrm{dvol}  $$
    ガンマ関数のおける大域的微分多様体は、これと同型により、
    $$ \int \Gamma \cdot {d \over {d \gamma}}{\Gamma dx_m} \le \int \Gamma
    dx_m + {d \over d \gamma}\Gamma $$
    $$ = \int \Gamma(\gamma)^{'}dx_m $$
    オイラーの定数の多様体積分は、
    $$ \int \left({{\int {1 \over x^s}}dx - \log x} \right)\mathrm{dvol} $$
    この解は、シュワルツシルト半径と同型より、
    $$ = e^f - e^{-f} \le {e^f + e^{-f}} $$
    次元の単位は多様体より、大域的微分多様体とオイラーの定数
    の多様体積分の加群分解は、オイラーの公式の三角関数の虚数の度解より、
    $$ {d \over df}F + \int C dx_m = 2 (\cos (i x \log x) 
    - i \sin (i x \log x)) $$
    すべては、大域的微分多様体の重力と反重力方程式に行き着き、
    $$ {d \over df}F \ge {d \over df}\int \int{1 \over (x \log x)^2}dx_m
    + {d \over df}\int \int{1 \over (y \log y)^{1 \over 2}}dy_m $$
    と求まる。

    これが、重力と反重力単独では、磁気単極子で、ド・ブロイのアイデアであり、
    両方では、磁気双極子で、南部・ゴールドストーンボソン粒子である。
    ヒッグス場の式でもある。

    北半球の磁気と南半球に磁気一セットで一種である。

    これらは、ハイゼンベルクとシュレーディンガー方程式のディラック作用素が、
    微分幾何からの作用素関数から、同一と言える。

    
    一般相対性理論における多様体積分とオイラーの定数の多様体積分
    が同型と言えることを上の式たちは述べている。

    このJones多項式が、金融市場の相場価格を決めるオプション方程式であり、
    流体力学による熱エントロピー値の熱エネルギー量が、
    いろんな機材に使われている。
    量子コンピュータにおける論理素子の回路の設計にも使われている。
    人工知能の論理素子と因数分解における量子暗号にも使われている。

    この一般相対性理論における多様体積分がJones多項式となり、
    金融市場で使われることを読んだのが、イスラエル国家らしい。

    $$ H\Psi = \bigoplus {\left(i \hbar^{\nabla} \right)}^{\oplus L} $$
    $$ = {1 \over 2}i e^{i \Hat{H}} $$
    $$ (i \hbar)^{'} = (- e^{i \Hat{H}})^{'} $$
    $$ = - i e^{i \Hat{H}} $$
    $$ \psi(x) = e^{-i \Hat{H}t}, \bigoplus (i \hbar ^{\nabla})^{\oplus L} 
    = {{{1 \over 2} e^{i \Hat{H}}}^{(-i e^{i \Hat{H}})}} $$
    $$ = ({1 \over 2}f)^{-i f}$$
    $$ = ({1 \over 2})^{-i f} \cdot e^{-x \log x} \cdot (f)^{i} $$
    $$ = \int e^{- x} x^{t - 1}dx, {d \over d \gamma}\Gamma = e^{-x \log x}  $$
     
    $$ f = x, i = t, {1 \over 2} = a, \bigoplus a^{-t x}x^t [I_m]
     \cong \int e^{-x} x^{t - 1}dx $$
     $$ |\psi(t)\rangle_s = e^{-i\Hat{H} t}|\Psi\rangle_{H},
    \Hat{A}_s = \Hat{A}_H(0) $$
    $$ |\Psi(t)\rangle_s \to {d \over dt} $$
    $$ i {d \over dt}|\psi (t) \rangle_s = \Hat{H}|\psi(t)\rangle_s $$
    $$ \langle\Hat{A}(t)\rangle = \langle\Psi(t)|\Hat{A}(0)|\Psi(t)\rangle  $$
    $$ {d \over dt}\Hat{A} = {1 \over i}[\Hat{A},H] $$
    $$ \Hat{A}(t) = e^{i\Hat{H}t}\Hat{A}(0)e^{-i\Hat{H}t} $$
    $$ \lim_{\theta \to 0} \begin{pmatrix} \sin \theta \\
    \cos \theta \end{pmatrix} \begin{pmatrix} \theta & 1 \\
    1 & \theta \end{pmatrix} \begin{pmatrix} \cos \theta \\
    \sin \theta \end{pmatrix} = \begin{pmatrix} 1 & 0 \\
    0 & - 1 \end{pmatrix} $$
    $$ f^{-1}(x) x f(x) = I^{'}_m, I^{'}_m = [1,0] \times [0,1] $$

    $$ {x + y} \ge \sqrt{x y} $$
    $$ {x^{{1 \over 2} + iy} \over e^{x \log x}} = 1 $$
    $$ \mathcal{O}(x) = \nabla_i \nabla_j \int e^{{2 \over m}\sin \theta 
    \cos \theta } \times {{N \mathrm{mod}(e^{x \log x})} \over {O(x)(x +
    \Delta |f|^2)^{1 \over 2}}} $$
    $$ x \Gamma(x) = 2 \int|\sin 2 \theta|^2 d \theta, \mathcal{O}(x)
    = m(x)[D^2 \psi] $$
    $$ i^2 = (0,1) \cdot (0,1), |a||b|\cos \theta = - 1 $$
    $$ E = \mathrm{div}(E, E_1) $$
    $$ \left({\{f,g\}} \over [f,g] \right) = i^2, E = mc^2, I^{'} = i^2 $$

    この式たちは、微分幾何の量子化から計算方法としての形式の
    大域的微分多様体の大域的部分積分が,大域的偏微分方程式としての式の
    形からわかるのも、微分幾何の量子化の仕組みとしての計算式から
    導けられるのも、すべては未来の私の子の情報から読み取れたと
    言える。
 
    $$ H \Psi = \bigoplus (i \hbar ^{\nabla})^{\oplus L} $$
    $$ = \bigoplus {{ H \Psi} \over \nabla L} $$
    $$ = e^{x \log x} = x^{(x)^{'}} $$
    however,
    $$ {d \over df}F = m(x) $$
    and gravity equation in being abled to existed with quantum physics,
    and this physics is also represented with quantum level of differential
    geometry.
    $$ {d \over dt}\psi(t) = \hbar $$
    $$ = {1 \over 2}i e^{i \Hat{H}} $$
    $$ (i \hbar)^{'} = (- e^{i \Hat{H}})^{'} $$
    $$ = - i e^{i \Hat{H}} $$
    $$ \psi(x) = e^{-i \Hat{H}t}, \bigoplus (i \hbar ^{\nabla})^{\oplus L} 
    = {{{1 \over 2} e^{i \Hat{H}}}^{(-i e^{i \Hat{H}})}} $$
    $$ = ({1 \over 2}f)^{-i f}$$
    $$ = ({1 \over 2})^{-i f} \cdot e^{-x \log x} \cdot (f)^{i} $$
    $$ = \int e^{- x} x^{t - 1}dx, {d \over d \gamma}\Gamma = e^{-x \log x}  $$
     
    $$ f = x, i = t, {1 \over 2} = a, \bigoplus a^{-t x}x^t [I_m]
     \cong \int e^{-x} x^{t - 1}dx $$
     Quantum level of differential geomerty is also resulted with
    Gamma function of global differential manifold of values, Heisenberg
    equation is alos resulted with global integrate manifold and this
    quantum level of differential geometry is manifold of deprive of formula.
    $$ |\psi(t)\rangle_s = e^{-i\Hat{H} t}|\Psi\rangle_{H},
    \Hat{A}_s = \Hat{A}_H(0) $$
    $$ |\Psi(t)\rangle_s \to {d \over dt} $$
    $$ i {d \over dt}|\psi (t) \rangle_s = \Hat{H}|\psi(t)\rangle_s $$
    $$ \langle\Hat{A}(t)\rangle = \langle\Psi(t)|\Hat{A}(0)|\Psi(t)\rangle  $$
    $$ {d \over dt}\Hat{A} = {1 \over i}[\Hat{A},H] $$
    $$ \Hat{A}(t) = e^{i\Hat{H}t}\Hat{A}(0)e^{-i\Hat{H}t} $$
    $$ \lim_{\theta \to 0} \begin{pmatrix} \sin \theta \\
    \cos \theta \end{pmatrix} \begin{pmatrix} \theta & 1 \\
    1 & \theta \end{pmatrix} \begin{pmatrix} \cos \theta \\
    \sin \theta \end{pmatrix} = \begin{pmatrix} 1 & 0 \\
    0 & - 1 \end{pmatrix} $$
    $$ f^{-1}(x) x f(x) = I^{'}_m, I^{'}_m = [1,0] \times [0,1] $$

    $$ {x + y} \ge \sqrt{x y} $$
    $$ {x^{{1 \over 2} + iy} \over e^{x \log x}} = 1 $$
    $$ \mathcal{O}(x) = \nabla_i \nabla_j \int e^{{2 \over m}\sin \theta 
    \cos \theta } \times {{N \mathrm{mod}(e^{x \log x})} \over {O(x)(x +
    \Delta |f|^2)^{1 \over 2}}} $$
    $$ x \Gamma(x) = 2 \int|\sin 2 \theta|^2 d \theta, \mathcal{O}(x)
    = m(x)[D^2 \psi] $$
    $$ i^2 = (0,1) \cdot (0,1), |a||b|\cos \theta = - 1 $$
    $$ E = \mathrm{div}(E, E_1) $$
    $$ \left({\{f,g\}} \over [f,g] \right) = i^2, E = mc^2, I^{'} = i^2 $$
     Gamma function of global differential equation is first of differential
    function deprivate with ordinary differential equations, more also
    universe and other dimensiion of Higgs field emerge with zeta function,
    more spectrum focus is three dimension of manifold of singularity or pair
    theorem, and universe and other dimension of each value equals.
    This reverse of circle function of manifolds is also Higgs fields of 
    dense of entropy and result equation. And differential structure of quantum
    level in gravity is poled with zeta function and quantum group in high 
    result values. Zeta funtion and quantum group is Global differential
    manifold of result in values.
    $$ {d \over d \gamma} \Gamma = m(x) $$
    $$ = e^{- x \log x} $$
    $$ \sin i x = {{e^{-x} + e^{x}} \over 2i} $$
    $$ {d \over df}F = m(x) = e^{f} + e^{-f} $$
    $$ = 2 i \sin (i {x \log x}) $$


   Beta function is,
      $$ \beta(p,q) = \int x^{1 - t} (1 - x)^{t}dx = \int t^{x}(1 - t)^{x - 1}
      dt $$
      $$ 0 \le y \le 1, \int^{1}_{0} x^{10}(1 - x)^{20}dx = B(11,21) $$
      $$ = {{\Gamma(11)\Gamma(21)} \over \Gamma(32)} = {{10! 20!} \over {
	      31!}} = {1 \over {931395465}} $$
       $$ {1 \over {931395465}} \cong {1 \over 9} = {1 \over {1 - x}} $$
	 $$     = {1 \over {1 - z}} =
	      \sum_{k=0}^{\infty}z^k = {1 \over {1 + z^2}} = \sum_{
		      k=0}^{\infty} (-1)^k z^{2k} $$
      $$ f(x) = \sum_{k=0}^{\infty} a_k z^k $$
      $$ {{d^n y} \over dx^n} = n! y^{n+1} $$
      $$ f^{(0)}(0) = n!f(0)^{n+1} = n! $$
      $$ f(x) \cong \sum_{k=0}^{\infty}x^n = {1 \over {1 - x}} $$
      In example script is,
      $$ {{dy} \over dx} = y^2, {1 \over y^2}\cdot {dy \over dx} = 1 $$
      $$ \int {1 \over y^2}{dy \over dx}dx = \int{dy \over y^2} = - {1 \over y
      } $$
      $$ - {1 \over y} = x - C, y = {1 \over {C - x}} $$
      $$ \exists x = 0, y = 1 $$
      $$\text{in first value condition compute with} $$
      $$ \text{result, C consumer sartified,} $$
      $$ y = {1 \over {1 - x}} $$
      This value result is concluded with native function from Abel manifold.

      Abelian enterstane with native funtion meltrage from Daranvelcian
      equation on Beta function.
 

   $$ H\Psi = \bigoplus {H \Psi \over \nabla L} $$
  　
 $$ \Gamma = \int e^{-x}x^{1 - t}dx $$
   $$ \gamma = \int e^{-x}x^{1 - t}\log x dx $$
   $$ = \int \Gamma(\gamma)^{'} dx_m$$
   $$ = {d \over d \gamma}\Gamma $$

   Eight differential geometry are each intersect with own level of
   concept from expalanation of Euler product system.
   This component of three manfold sergeried with geometry of destroy and
   desect with time element of Stokes equation.
   ８種類の微分幾何では、それぞれの時間の固有値が違う。
   閉3次元多様体上で微分幾何の切除、分解によって時間の性質が決まる。
   This manifold gut theory from described with zeta function to
   catastrophe for non tree of routs result on sergery of space system.
   閉３次元多様体に統合されると、ゼータ関数になり、各幾何に分解された場合
   に、非分岐から、この上、sergeryの結果が決まる。

   各微分構造においての時間発展での熱エネルギーの変化は$E=mc^2, E=m^2 c^2 $
   のthree manifoldが微分幾何構造の変化にそれぞれ対応する。
   Each geometry point interacte with exchange of three manifold in
   Seifert structure from special relativity of equation for
   heat energy fluentations on time developyed from this extermaination.

   $$ \nabla ({i \hbar^{\nabla}})^{\oplus L}, \bigoplus ({i \hbar^{\nabla}})^{
	   \oplus L}, \square ({i \hbar^{\nabla}})^{\oplus L} $$
   These operator of equation on summatative manifolds from emerged with
   element of particle conclution, Euler product is resulted from this
   operator expalanation of locality insectations.
   これらの作用素が加減乗除の生成式の元、Euler Productの結論による
   作用素生成の論理素子でもある。
   Moreover, this eight geometry of differential operator are
   constructed with four pattern of Jones manifold from
   summatative formula and this system extate with special relativity
   references. This decieved of elemet on summatative equation routed
   to internext in real and imaginary pole on complex dimension,
p  and this dimension explation with Stokes theorem defined too.
   And this defined circutation of Yacobi matrix is represected with
   Knot theory with anstate with Jones manifold.
   その上に、8種類の微分幾何は、Jones多項式の4パターンでの構成される
   差分と加群方程式から、特殊相対性理論をも思わせる、この差分方程式
   $$ {d \over d \gamma}\Gamma = e^{f} + e^{-f} \ge e^{f} - e^{f} $$
   から、8種類の代数幾何が、ストークスの定理と同じく、このJones多項式
   の固有周期が、すべての時間の流れの基軸をも内包してもいることが、
   わかる。このJones多項式は、結び目の理論をも、この周期が言い表してもいる。
   $$ {d \over d \gamma}\Gamma + \int C dx_m = 2(\cos(i x \log x) - i \sin
   (i x \log x)) $$
    $$ e^{- \theta} = \cos(i x \log x) - i \sin(i x \log x)$$
    $$ \int \Gamma(\gamma)^{'}dx_m = 2(\cos (i x \log x) 
    - i \sin(i x \log x)) $$

 カルーツァ・クライン空間の方程式は、
    $$ ds^2 = g_{\mu\nu}(x)dx^{\mu}dx^{\nu} + \psi^2(x)(dx^2 + \kappa^2
    A_{\mu}(x)dx^{\nu})^2 $$
    この式は、
    $$ ds^2 = -N(r)^2dt^2 + \psi^2(x)(dr^2 + r^2d\theta^2) $$
    $$ ds^2 = -dt^2 + r^{-8\pi Gm}(dr^2 + r^2d\theta^2) $$
    とまとまり、
    $$ dx^2 = g_{\mu\nu}(x)(g_{\mu\nu}(x)dx^2 - dxg_{\mu\nu}(x)) $$
    $$ dx = (g_{\mu\nu}(x)^2dx^2 - g_{\mu\nu}(x)dx g_{\mu\nu}(x))^{1 \over 2} $$
    と反重力と正規部分群の経路和となり、　
    $$ \pi(\chi, x) = i \pi (\chi, x)f(x) - f(x)  \pi (\chi, x) $$
    と基本群にまとまる。
    シュワルツシルト半径は、
    $$ ds^2 = -(1 - {r_s \over r})c^2 dt^2 + {1 \over {{1 - {r_s \over r}}}} 
    dr^2 + r^2 d \theta^2 + r^2 \sin^2 \theta d \psi^2 $$
    これは、カルーツァ・クライン空間と同型となる。
    一般相対性理論は、
    $$ R^{\mu\nu} + {1 \over 2}g_{ij}\Lambda = \kappa T^{\mu\nu} $$
    この多様体積分は、
    $$ \int \kappa T^{\mu\nu}\mathrm{dvol} = \int  \left(R^{\mu\nu} + 
    {1 \over 2}g_{ij}\Lambda\right)\mathrm{dvol}  $$
    ガンマ関数のおける大域的微分多様体は、これと同型により、
    $$ \int \Gamma \cdot {d \over {d \gamma}}{\Gamma dx_m} \le \int \Gamma
    dx_m + {d \over d \gamma}\Gamma $$
    $$ = \int \Gamma(\gamma)^{'}dx_m $$
    オイラーの定数の多様体積分は、
    $$ \int \left({{\int {1 \over x^s}}dx - \log x} \right)\mathrm{dvol} $$
    この解は、シュワルツシルト半径と同型より、
    $$ = e^f - e^{-f} \le {e^f + e^{-f}} $$
    次元の単位は多様体より、大域的微分多様体とオイラーの定数
    の多様体積分の加群分解は、オイラーの公式の三角関数の虚数の度解より、
    $$ {d \over df}F + \int C dx_m = 2 (\cos (i x \log x) 
    - i \sin (i x \log x)) $$
    すべては、大域的微分多様体の重力と反重力方程式に行き着き、
    $$ {d \over df}F \ge {d \over df}\int \int{1 \over (x \log x)^2}dx_m
    + {d \over df}\int \int{1 \over (y \log y)^{1 \over 2}}dy_m $$
    と求まる。



       超重力理論が、超弦理論と一般相対性理論の多様体積分へと、分岐している。
    この超弦理論と一般相対性理論の多様体積分が、超重力理論へと統合される。
    
    $\theta = 45$度は、平行四辺形の、一般座標空間としての変形をすると、
    $ \theta = 60$度へと、線形変換することで、$\sin \theta = {\sqrt{3} \over
    2} $
    となる。循環となると、永遠の寿命もだめになり、木星の大善は、
    難しいのは、一目瞭然であり、仏教の前後関係もある。

\end{document}
